% Basic catcodes
\catcode`\{=1 \catcode`\}=2 \catcode`\$=3 \catcode`\&=4 \catcode`\#=6 \catcode`\\=0 \catcode`\^=7 \catcode`\_=8
\catcode`\ =10 \catcode`\^^M=5 \catcode`\a=11 \catcode`\z=11 \catcode`\A=11 \catcode`\Z=11

\directlua{ tex.enableprimitives('', tex.extraprimitives()) }

% Define the font using the integrated loader.
\font\tenmath="[FiraMath-Regular.otf]" at 10pt

% Assign the font to the math families
\textfont0=\tenmath
\scriptfont0=\tenmath
\scriptscriptfont0=\tenmath

% Verification Block
\directlua{
  local f = font.getfont(font.current())
  if f then
    if f.MathConstants then
      print("SUCCESS: 'MathConstants' table found.")
      % Example: Print a specific value to verify (e.g., AxisHeight)
      % Values are usually in font design units
      print("AxisHeight:", f.MathConstants.AxisHeight)
      
      % Force a parameter update if the engine didn't do it automatically
      % (Standard LuaTeX does this automatically when MathConstants is present)
      if tex.getmath("mathaxis", "display") == 0 then
          print("WARNING: Math parameters were not applied automatically.")
       else
          print("CONFIRMED: Engine applied MathConstants to internal parameters.")
      end
    else
      print("ERROR: No 'MathConstants' table found. Check font loader.")
    end
  end
}

% Standard mapping for the test
\def\loop#1\repeat{\def\body{#1}\iterate}
\def\iterate{\body \let\next\iterate \else\let\next\relax\fi \next}
\count0=`0 \loop \Umathcode\count0 = 7 0 \count0 \ifnum\count0 < `9 \advance\count0 by 1 \repeat
\count0=`a \loop \Umathcode\count0 = 7 0 \count0 \ifnum\count0 < `z \advance\count0 by 1 \repeat
\count0=`A \loop \Umathcode\count0 = 7 0 \count0 \ifnum\count0 < `Z \advance\count0 by 1 \repeat

\Umathcode`+=2 0 `+
\Umathcode`-=2 0 `-
\Umathcode`(=4 0 `(
\Umathcode`)=5 0 `)
\Umathcode`/=0 0 `/
\Umathcode`*=2 0 `*

% Delimiter codes
% Note: We use `\{ and `\} on the left because backtick-escape works for catcode 1/2.
% But on the right (the value), we must strictly provide the slot number.

\Udelcode`( = 0 0 `(
\Udelcode`) = 0 0 `)
\Udelcode`[ = 0 0 `[
\Udelcode`] = 0 0 `]

% For braces, use the numeric ASCII values 123 ("7B) and 125 ("7D)
% Simplest safe way:
\Udelcode`\{ = 0 0 "7B   % Hex 7B is 123 ({)
\Udelcode`\} = 0 0 "7D   % Hex 7D is 125 (})

\Udelcode`| = 0 0 `|

% Initialize the main text font to something usable
\font\mainfont="[FiraMath-Regular.otf]" at 10pt
\mainfont
\def\bye{\par\vfill\supereject\end}
\def\fmtname{plain} % Helpful for some loggers

% \over is a primitive in TeX, so \frac is just a wrapper
\def\frac#1#2{{#1\over#2}}

% \sqrt uses the \Uradical primitive for OpenType math
% Syntax: \Uradical <fam> <char>
% "221A is the Unicode slot for the square root symbol
\def\sqrt{\Uradical 0 "221A }

\dump

